\section{Testing stationary approximations to edge flow}
\label{sec:trial-flow}
For a stationary set $F$ the exact probability of returning in one step,
conditional on starting in $F$, is $J_h(F,F)/\nu(F)$, where
\[
 J_v(x,y)=\mu(x)P(x,y)v(x)v(y),\qquad \nu_v(x)=\mu(x)v(x)Pv(x).
\]
Here $\nu=\nu_h$. A candidate computable from $\pi$ alone is $u=\pi/(P\sqrt\pi)$.
We give two residual tests that apply to any positive trial $u$. We also test
this particular choice against the programme's matched Hamming shell.

\paragraph{Two stationary thresholds at the critical moment.}
In the setting of (SR1) one can further restrict its set $Z$ to
$D=Z\cap\{\pi>e^{\beta t}\}$ for any fixed $0<\beta<4/9$. Indeed
Section~\ref{sec:matched-defect} gives $\|1/g\|_6\le C_B$.
Thus $\|U1_{\{\pi\le e^{\beta t}\}}\|_3\le C_B^2 e^{\beta t}$.
On $A=\{U>e^t\}$ truncating from order $3$ to $5/3$ gives a bound
$C'_B\exp[-(4-9\beta)t/5]$. The portion of $Z$ off $A$ was already
estimated in (SR3). Add the two nonnegative weights between the two
copies of $Q$ as in (SR1). This proves $|s_Z-s_D|\le Ce^{-ck}$, hence
full ordered recovery is equivalent to $R_D\to0$. The rest of this
section concerns how to evaluate returns once a set has been selected.

\paragraph{An entropy characterization of nonnegative spectrum.}
Let $P$ now be any strictly positive finite Markov kernel reversible
with respect to $\mu$. For positive measures use the unnormalized
relative entropy $H(A\mid B)=\sum[A\log(A/B)-A+B]$. The following
property is equivalent to $P$ being positive semidefinite in $L^2(\mu)$:
for every positive $h,u$ with $\langle h,Ph\rangle_\mu=
\langle u,Pu\rangle_\mu=1$,
\[
       H(\nu_h\mid\nu_u)\le H(J_h\mid J_u)
                    \le 2H(\nu_h\mid\nu_u).                    \tag{EF1}
\]
For positive semidefinite $P$ both inequalities continue to hold for
unnormalized trials. The constant two is attained even for normalized
inputs when $P(x,y)=\mu(y)$ and the two induced marginals differ.
Indeed the normalized edge laws then factor as products of their marginals.

The first inequality is marginalization. For the second put $a=u/h$,
$K(x,y)=P(x,y)h(y)/Ph(x)$, and $\phi(z)=z-1-\log z$.
Expanding the logarithms and using reversibility gives the exact identity
\[
  2H(\nu_h\mid\nu_u)-H(J_h\mid J_u)=
  \langle u-h,P(u-h)\rangle_\mu+2\sum_x\nu_h(x)\phi(Ka(x)).  \tag{EF2}
\]
In particular the right side is nonnegative in the claimed regime.
To prove necessity, suppose that $P f=\lambda f$ with $\lambda<0$,
$\mathbb E_\mu f=0$ and $\|f\|_2=1$. Strict positivity ensures $\lambda>-1$.
Take $h=1$ and $u_s=(1+sf)/\sqrt{1+\lambda s^2}$ for sufficiently small
real $s$. The scores at $s=0$ of $J_{u_s}$ and $\nu_{u_s}$ are
$f(x)+f(y)$ and $(1+\lambda)f(x)$ respectively. A normalized finite
family with score $b$ has relative entropy from its initial law
$s^2\mathbb E b^2/2+O(s^3)$ in either orientation here, by Taylor
expansion. The ratio of edge to marginal entropy tends to
$2/(1+\lambda)>2$, disproving the upper inequality.

\paragraph{Optimal factor for a fixed kernel.}
Assume at least two states and strict positivity as above. Define
\[
 \kappa(P)=\sup_{\substack{f,g\ge0,\ f,g\ne0\\ f g=0}}
 \frac{\langle f,Pg\rangle_\mu}
 {\sqrt{\langle f,Pf\rangle_\mu\langle g,Pg\rangle_\mu}} .
\]
Both denominators are positive by the diagonal. The least constant
$C(P)$ in the normalized upper inequality of (EF1), with $2$ replaced
by $C(P)$, is
\[
                   C(P)=\max\{2,1+\kappa(P)\}.               \tag{EF5}
\]
In particular a strictly positive definite kernel still has best
global constant two. Highly uneven positive scalings are allowed.

For a probability density $v$ with positive scaling $h_v$, set
$G(v)=\frac12 H(J_{h_v}\mid\mu P)$, where $\mu P(x,y)=\mu(x)P(x,y)$.
Since $G(v)=\langle v,\log h_v\rangle_\mu$, its derivative along
the probability simplex is $\log h_v$: the remaining differential
$\langle v,d\log h_v\rangle_\mu$ is half the variation of the unit
edge mass. Thus the edge entropy from $v$ to $w$ is twice
$G(v)-G(w)-\langle v-w,\log h_w\rangle_\mu$.
The Hessian of $G$ on a tangent of score $b=dv/v$ is
$\langle b,(I+K_v)^{-1}b\rangle_{v\mu}$, by logarithmic
differentiation of $h_v Ph_v=v$. Existence and smoothness here use the
positive diagonal and strict positivity of $P$; positive semidefiniteness
is unnecessary. The entropy of the marginal has Hessian
$\|b\|_{L^2(v\mu)}^2$. Integrating both Hessians with weight $1-z$
along the line from $w$ to $v$, and taking arbitrarily short lines
in the least-eigenvalue direction for the converse, gives
$C(P)=2/(1+\inf_v\lambda_{\min}(K_v))$. The constant eigenfunction has
eigenvalue one and all tangents have mean zero.

To evaluate the infimum write $\xi=h f_0$ for a test function $f_0$
in the Rayleigh quotient of $K_{hPh}$. Its numerator is
$\langle\xi,P\xi\rangle_\mu$ and its denominator is
$\sum_{x,y}\mu(x)P(x,y)\xi(x)^2 h(y)/h(x)$.
The latter is at least $\langle|\xi|,P|\xi|\rangle_\mu$ by pairing
the two orientations and applying arithmetic--geometric mean.
Taking $h=|\xi|+\epsilon$ attains equality in the limit, up to the
irrelevant scale of $h$. For a negative numerator this gives the
most negative quotient. Write $\xi=f-g$ with nonnegative disjoint
parts, and put $A=\langle f,Pf\rangle_\mu$, $B=\langle g,Pg\rangle_\mu$,
$C=\langle f,Pg\rangle_\mu$. The ratio for the optimal denominator is
$(A+B-2C)/(A+B+2C)$. Rescaling one part to make $A=B$ minimizes it
for fixed shapes. Hence a negative infimum exists precisely when
$\kappa(P)>1$, and then equals $(1-\kappa(P))/(1+\kappa(P))$.
In the remaining case every numerator is nonnegative. Keeping $h$ fixed
at one on a single state and tending to zero elsewhere makes all rows
of $K_{hPh}$ tend to the point mass on that state. The nonconstant
eigenvalues then tend to zero by continuity of the characteristic
polynomial. This proves (EF5).

\paragraph{A computable error bound for a prescribed return.}
For Boolean noise, calculate $H_m=H(\nu\mid\nu_u)$ from the target and
the ordinary kernel. For any nonempty $F$ the log-sum inequality on
$F\times F$, followed by $(\sqrt v-\sqrt w)^2\le v\log(v/w)-v+w$,
gives
\[
 \left|\sqrt{J_h(F,F)/\nu(F)}-\sqrt{J_u(F,F)/\nu(F)}\right|
                 \le\sqrt{2H_m/\nu(F)}.                      \tag{EF3}
\]
In particular $H_m=o(m_0)$ suffices simultaneously for every tested
set with $\nu(F)\ge m_0$. This bound uses the actual stationary mass
even for an unnormalized trial. A small absolute residual alone can
be inconclusive on exponentially rare sets.

\paragraph{A logarithmic mass cost under uniform relative control.}
Normalize the trial to have edge mass one and set
$\eta=\|\log(\nu_u/\nu)\|_\infty$. For any nonempty edge event $E$ put
$s_0=J_h(E)$ and $s_1=J_u(E)$. There is a numerical $C$ such that
\[
 e^{-C\eta}\le
 \frac{1+\log(1/s_1)}{1+\log(1/s_0)}
 \le e^{C\eta}.                                                \tag{EF4}
\]
Here $P$ is strictly positive, reversible and positive semidefinite.
To prove this, interpolate the marginal probabilities as
$\nu_z\propto\nu^{1-z}\nu_u^z$ for $0\le z\le1$.
Write $h_z$ for the exact scaling along the path. Existence, uniqueness
and smoothness follow by minimizing the symmetric exponential scaling
functional used in the \href{https://airr.science/papers/ARR-2026-2B2W1WQ7DF9N5VW7/}{matched-shell proof}:
its diagonal terms ensure coercivity and a positive definite Hessian,
and the linear term uses the desired positive marginal. Thus
$(I+K_z)\partial_z\log h_z=b_z=\partial_z\log\nu_z$,
with $\|b_z\|_\infty\le2\eta$. The reversible $K_z$ has spectrum in
$[0,1]$ by diagonal congruence and stochasticity. Put $T_z=(I-K_z)/2$.
Its infinity norm is at most one and its $L^2(\nu_z)$ norm at most $1/2$.
Use $(I+K_z)^{-1}=\frac12\sum_{j\ge0}T_z^j$. Conditional on the event
$E$ of probability $s=J_{h_z}(E)$, each endpoint has density at most
$1/s$ relative to $\nu_z$ and hence $L^2$ norm at most $s^{-1/2}$.
Differentiating $\log s$ averages $\partial_z\log h_z$ at both endpoints.
Bounding each term in the series in both norms yields
\[
 |\partial_z\log s|\le 2\eta\sum_{j\ge0}\min(1,2^{-j}s^{-1/2})
                  \le C\eta(1+\log(1/s)).
\]
Integration proves (EF4). If $\eta(1+\log(1/m_0))$ stays bounded, a
conditional return bounded below by a fixed positive number for a set
of mass at least $m_0$ under either edge law (using $\nu(F)$ as the
common denominator) is bounded below under the other. Indeed start at
the endpoint with $s_i\ge c\,\nu(F)$ and apply (EF4) in the other
direction: the change in $\log s$ is at most
$(e^{C\eta}-1)(1+\log(1/s_i))=O_c(1)$. If the product tends to zero,
the multiplicative error for such persistent returns tends to one.
This certificate uses the trial's marginal residual and a logarithm
of the tested mass. No estimate of the mixing gap enters.

\paragraph{A false negative for the one-step formula on a matched shell.}
Use the positive-floor, single-basis half-noise construction of the
\href{https://airr.science/papers/ARR-2026-2B2W1WQ7DF9N5VW7/versions/v1/}{matched Hamming shell}.
For its target sphere $S$ let $q=P(x,S)$ for $x\in S$, and let
$\pi_S$ be the constant stationary density there. In the source proof,
with a fixed small exponent $\kappa>0$, the spike amplitude is
$M=e^{-\kappa L/2}/q$, its potential on $S$ is at most $2M$ before
normalization, and its noise average there is at most $3e^{-\kappa L/2}$.
The normalizing constant tends to one. Exact marginal correction changes
the density by $O(e^{-cL})$ uniformly. Thus $q\pi_S=O(e^{-c'L})$ while
the actual conditional return to $S$ tends to one exponentially.
For $u=\pi/(P\sqrt\pi)$, the floor $\pi\ge b>0$ gives $u\le\pi/\sqrt b$.
For every nonempty $F\subseteq S$,
\[
        J_u(F,F)/\nu(F)\le q\pi_S/b=O(e^{-c'L}).
\]
Even normalizing the trial edge law does not repair this false negative.
The shell construction has $\|\pi-1\|_3=O(e^{-cL})$. Hence
$\|P\sqrt\pi-1\|_2\le\|\pi-1\|_2$ and $\|u-1\|_2=O_b(e^{-cL})$.
Consequently the trial edge mass tends to one, and its marginal density
$uPu$ tends to $\pi$ exponentially in $L^1$ by Cauchy--Schwarz.
The example concerns the actual returning sphere in the single-basis
model. A floor makes $Z$ empty for large $k$; this example does not
decide the residual-free trial formula on the two-base critical set $D$.

\paragraph{Relation to scaling theory.}
Symmetric scaling algorithms such as those of
\href{https://strathprints.strath.ac.uk/49349/1/main_revised3.pdf}{Knight, Ruiz and U\c{c}ar}
and entropy estimates for entropic plans such as those of
\href{https://www.math.columbia.edu/~mnutz/docs/Sinkhorn_rate.pdf}{Ghosal and Nutz}
provide relevant comparisons. A closer antecedent for the normalized
upper inequality in (EF1) is the convexity of the Sinkhorn negentropy
in href{https://proceedings.mlr.press/v89/feydy19a.html}{Feydy et al.,
Propositions 3--4}. For the kernel density $P(x,y)/mu(y)$ and unit
regularization, their self-transport functional with the negative
one-half convention equals marginal entropy relative to $mu$ minus
$G$, up to an affine term. Its convexity gives the normalized upper
inequality directly under their positive-kernel assumptions. Formula
(EF5) also treats kernels with a negative eigenvalue. Here (EF2) concerns two symmetric diagonal
scalings of the same positive kernel and uses its quadratic form directly.
The event estimate (EF4) weights the inverse scaling derivative by the
two endpoint laws conditioned on an event.
